\documentclass[12pt,a4paper]{ctexbook}
\usepackage{geometry}
\geometry{margin=2.2cm}
\usepackage{setspace}
\setstretch{1.35}
\usepackage{hyperref}
\title{全宋词选·ci.song.8000}
\author{古典文献整理}
\date{}
\begin{document}
\maketitle
\tableofcontents
\newpage
\section{采莲实催・采莲令}
\textbf{史浩}\par\vspace{0.5em}
露华霞液，云桨椒醑，姿玉斝金罍。\par
交酬成雅会。\par
拚沈醉。\par
中山千日，未为长久，今此陶陶一饮，动经万祀。\par
陈果蓏，皆是奇异。\par
似瓜如斗尽备。\par
三千岁。\par
一熟珍味。\par
饤坐中，莹似玉、爽口流涎，三偷不枉，西真指议。\par

\section{采莲衮・采莲令}
\textbf{史浩}\par\vspace{0.5em}
有珍馔，时时馈。\par
滑甘丰腻。\par
紫芝荧煌，嫩菊秀媚。\par
贮玛瑙琥珀精器。\par
延年益寿莫拟。\par
人间烹饪徒费。\par
休说龙肝凤髓。\par
动妙乐、仙音鼎沸。\par
玉箫清，瑶瑟美。\par
龙笛脆。\par
杂还飞鸾，花裀上、趁拍红牙，馀韵悠扬，竟海变桑田未止。\par

\section{采莲歇拍・采莲令}
\textbf{史浩}\par\vspace{0.5em}
其间有洞天侣，思游尘世。\par
珠葆摇曳。\par
华表真人，清江使者，相从密议。\par
此老遨嬉。\par
我辈应须随侍。\par
正举步、忽思同类。\par
十八公、方耸壑，宜邀致。\par
夙驾星言，人争图绘。\par
朅来鄞山甬水。\par
因此崇成，四明里第。\par

\section{采莲煞衮・采莲令}
\textbf{史浩}\par\vspace{0.5em}
吾皇喜。\par
光宠无贰。\par
玉带金鱼荣贵。\par
或者疑之。\par
岂识圣明，曾主斯乡，尝相与尽缱绻，胶漆何可相离。\par
今日风云合契。\par
此实天意。\par
吾皇圣寿无极，享晏粲千载相逢，我翁亦昌炽。\par
永作升平上瑞。\par

\section{采莲舞}
\textbf{史浩}\par\vspace{0.5em}
练光浮，烟敛澄波渺。\par
燕脂湿、靓妆初了。\par
绿云伞上露滚滚，的皪真珠小。\par
笼娇媚、轻盈伫眺。\par
无言不见仙娥，凝望蓬岛。\par
玉阙葱葱，镇锁佳丽春难老。\par
银潢急、星槎飞到。\par
暂离金砌，为爱此、极目香红绕。\par
倚兰棹。\par
清歌缥缈。\par
隔花初见，楚楚风流年少。\par

\section{采莲舞}
\textbf{史浩}\par\vspace{0.5em}
蕊沼清冷涓滴水。\par
迢迢烟浪三千里。\par
微孕青房包绣绮。\par
薰风里。\par
幽芳洗尽闲桃李。\par
羽氅飘萧尘外侣。\par
相呼短棹轻偎倚。\par
一片清歌天际起。\par
声尤美。\par
双双惊起鸳鸯睡。\par

\section{采莲舞}
\textbf{史浩}\par\vspace{0.5em}
翠盖参差森玉柄。\par
迎风浥露香无定。\par
不著尘沙真体净。\par
芦花径。\par
酒侵酥脸霞相映。\par
棹拨木兰烟水暝。\par
月华如练秋空静。\par
一曲悠扬沙鹭听。\par
牵清兴。\par
香红已满蒹葭艇。\par

\section{采莲舞}
\textbf{史浩}\par\vspace{0.5em}
草软沙平风掠岸。\par
青蓑一钓烟江畔。\par
荷叶为裀花作幔。\par
知谁伴。\par
醇醪只把鲈鱼换。\par
盘缕银丝杯自暖。\par
篷窗醉著无人唤。\par
逗得醒来横脆管。\par
清歌缓。\par
彩鸾飞去红云乱。\par

\section{采莲舞}
\textbf{史浩}\par\vspace{0.5em}
太华峰头冰玉沼。\par
开花十丈干云杪。\par
风散天香闻四表。\par
知多少。\par
亭亭碧叶何曾老。\par
试问霏烟登鸟道。\par
丹崖步步祥光绕。\par
折得一枝归月峤。\par
蓬莱岛。\par
霞裾侍女争言好。\par

\section{采莲舞}
\textbf{史浩}\par\vspace{0.5em}
珠露漙漙清玉宇。\par
霞标绰约消烦暑。\par
时驭清风之帝所。\par
寻旧侣。\par
三千仙仗临烟渚。\par
舴艋飘摇来复去。\par
渔翁问我居何处。\par
笑把红蕖呼鹤驭。\par
回头语。\par
壶中自有朝天路。\par

\section{采莲舞}
\textbf{史浩}\par\vspace{0.5em}
彤霞出水弄幽姿。\par
娉婷玉面相宜。\par
棹歌先得一枝枝。\par
波上画鲸飞。\par
向此画堂高会，幽馥散、堪引瑶卮。\par
幸然逢此太平时。\par
不醉可无归。\par

\section{采莲舞}
\textbf{史浩}\par\vspace{0.5em}
蕊宫阆苑。\par
听钧天帝乐，知他几遍。\par
争似人间，一曲采莲新传。\par
柳腰轻，莺舌啭。\par
逍遥烟浪谁羁绊。\par
无奈天阶，早已催班转。\par
却驾彩鸾，芙蓉斜盼。\par
愿年年，陪此宴。\par

\section{太清舞}
\textbf{史浩}\par\vspace{0.5em}
武陵自古神仙府。\par
有渔人迷路。\par
洞户迸寒泉，泛桃花容与。\par
寻花逶逦见灵光，舍扁舟、飘然入去。\par
注目渺红霞，有人家无数。\par

\section{太清舞}
\textbf{史浩}\par\vspace{0.5em}
须臾却有人相顾。\par
把肴桨来聚。\par
礼数既雍容，更衣冠淳古。\par
渔人方问此何乡，众颦眉、皆能深诉。\par
元是避嬴秦，共携家来住。\par

\section{太清舞}
\textbf{史浩}\par\vspace{0.5em}
当时脱得长城苦。\par
但熙熙朝暮。\par
上帝锡长生，任跳丸乌兔。\par
种桃千万已成阴，望家乡、杳然何处。\par
从此与凡人，隔云霄烟雨。\par

\section{太清舞}
\textbf{史浩}\par\vspace{0.5em}
渔舟之子来何所。\par
尽相猜相语。\par
夜宿玉堂空，见火轮飞舞。\par
凡心有虑尚依然，复归指、维舟沙浦。\par
回首已茫茫，叹愚迷不悟。\par

\section{太清舞}
\textbf{史浩}\par\vspace{0.5em}
我今来访烟霞侣。\par
沸华堂箫鼓。\par
疑是奏钧天，宴瑶池金母。\par
却将桃种散阶除，俾华实、须看三度。\par
方记古人言，信有缘相遇。\par

\section{太清舞}
\textbf{史浩}\par\vspace{0.5em}
云軿羽幰仙风举。\par
指丹霄烟雾。\par
行作玉京朝，趁两班鹓鹭。\par
玲珑环佩拥霓裳，却自有、箫韶随步。\par
含笑嘱芳筵，后会须来赴。\par

\section{太清舞}
\textbf{史浩}\par\vspace{0.5em}
游尘世、到仙乡。\par
喜君王。\par
跻治虞唐。\par
文德格遐荒。\par
四裔尽来王。\par
干戈偃息岁丰穰。\par
三万里农桑。\par
归去告穹苍。\par
锡圣寿无疆。\par

\section{柘枝舞}
\textbf{史浩}\par\vspace{0.5em}
□人奉圣□□朝□□□□主□□□□□留伊。\par
得荷云戏、幸遇文明，尧阶上、太平时。\par
□□□□何不罢岁□征舞柘枝。\par

\section{柘枝舞}
\textbf{史浩}\par\vspace{0.5em}
我是柘枝娇女。\par
□□多风措。\par
□□□、□住深□，妙学得柘枝舞。\par
□□头戴凤冠，□□□纤腰束素。\par
□□遍体锦衣装，来献呈歌舞。\par

\section{柘枝舞}
\textbf{史浩}\par\vspace{0.5em}
回头却望尘寰去。\par
喧画堂箫鼓。\par
整云鬟、摇曳青绡，爱一曲柘枝舞。\par
好趁华封盛祝笑，共指南山烟雾。\par
蟠桃仙酒醉升平，望凤楼归路。\par

\section{花舞}
\textbf{史浩}\par\vspace{0.5em}
贵客之名从此有。\par
多谢风流，飞驭陪尊酒。\par
持此一卮同劝后。\par
愿花长在人长寿。\par

\section{花舞}
\textbf{史浩}\par\vspace{0.5em}
嘉客之名从此有。\par
多谢风流，飞驭陪尊酒。\par
持此一卮同劝后。\par
愿花长在人长寿。\par

\section{花舞}
\textbf{史浩}\par\vspace{0.5em}
素客之名从此有。\par
多谢风流，飞驭陪尊酒。\par
持此一卮同劝后。\par
愿花长在人长寿。\par

\section{花舞}
\textbf{史浩}\par\vspace{0.5em}
幽客之名从此有。\par
多谢风流，飞驭陪尊酒。\par
持此一卮同劝后。\par
愿花长在人长寿。\par

\section{花舞}
\textbf{史浩}\par\vspace{0.5em}
野客之名从此有。\par
多谢风流，飞驭陪尊酒。\par
持此一卮同劝后。\par
愿花长在人长寿。\par

\section{花舞}
\textbf{史浩}\par\vspace{0.5em}
雅客之名从此有。\par
多谢风流，飞驭陪尊酒。\par
持此一卮同劝后。\par
愿花长在人长寿。\par

\section{花舞}
\textbf{史浩}\par\vspace{0.5em}
净客之名从此有。\par
多谢风流，飞驭陪尊酒。\par
持此一卮同劝后。\par
愿花长在人长寿。\par

\section{花舞}
\textbf{史浩}\par\vspace{0.5em}
仙客之名从此有。\par
多谢风流，飞驭陪尊酒。\par
持此一卮同劝后。\par
愿花长在人长寿。\par

\section{花舞}
\textbf{史浩}\par\vspace{0.5em}
寿客之名从此有。\par
多谢风流，飞驭陪尊酒。\par
持此一卮同劝后。\par
愿花长在人长寿。\par

\section{花舞}
\textbf{史浩}\par\vspace{0.5em}
清客之名从此有。\par
多谢风流，飞驭陪尊酒。\par
持此一卮同劝后。\par
愿花长在人长寿。\par

\section{花舞}
\textbf{史浩}\par\vspace{0.5em}
近客之名从此有。\par
多谢风流，飞驭陪尊酒。\par
持此一卮同劝后。\par
愿花长在人长寿。\par

\section{花舞}
\textbf{史浩}\par\vspace{0.5em}
算仙家，真巧数，能使众芳长绣组。\par
羽軿芝葆，曾到世间，谁共凡花为伍。\par
桃李漫夸艳阳，百卉又无香可取。\par
岁岁年年长是春，何用芳菲分四序。\par

\section{花舞}
\textbf{史浩}\par\vspace{0.5em}
对芳辰，成良聚，珠服龙妆环宴俎。\par
我御清风，来此纵观，还须折枝归去。\par
归去蕊珠绕头，一一是东君为主。\par
隐隐青冥怯路遥，且向台中寻伴侣。\par

\section{花舞}
\textbf{史浩}\par\vspace{0.5em}
叹尘寰，乌兔走，花谢花开能几许。\par
十分春色，一半遣愁，那堪飘零风雨。\par
争似此花自然，悄不待、根生下土。\par
花既无凋春又长，好带花枝倾寿醑。\par

\section{花舞}
\textbf{史浩}\par\vspace{0.5em}
是非场，名利海，得丧炎凉徒自苦。\par
至乐陶陶，唯有醉乡，谁向此间知趣。\par
花下一杯一杯，且莫把、光阴虚度。\par
八极神游长寿仙，蜾蠃螟蠕休更觑。\par

\section{剑舞}
\textbf{史浩}\par\vspace{0.5em}
荧荧巨阙。\par
左右凝霜雪。\par
且向玉阶掀舞，终当有、用时节。\par
唱彻。\par
人尽说。\par
宝此制无折。\par
内使奸雄落胆，外须遣、豺狼灭。\par

\section{渔夫舞}
\textbf{史浩}\par\vspace{0.5em}
细雨斜风都不管。\par
柔蓝软绿烟堤畔。\par
鸥鹭忘机为主伴。\par
无羁绊。\par
等闲莫许金章换。\par

\section{渔夫舞}
\textbf{史浩}\par\vspace{0.5em}
好景侬家披得去。\par
前村雪屋云深处。\par
一棹清歌归晚浦。\par
真佳趣。\par
知谁画得归缣素。\par

\section{渔夫舞}
\textbf{史浩}\par\vspace{0.5em}
济涉还渠渔父子。\par
生涯只在烟波里。\par
练静忽然风又起。\par
赢得底。\par
吹来别浦看桃李。\par

\section{渔夫舞}
\textbf{史浩}\par\vspace{0.5em}
万斛黄金迷俯仰。\par
轻舠不碍飞双桨。\par
光透碧霄千万丈。\par
真堪赏。\par
恰如镜里人来往。\par

\section{渔夫舞}
\textbf{史浩}\par\vspace{0.5em}
已见白鱼翻翠荇。\par
任公一掷波千顷。\par
不是六鳌休便领。\par
清昼永。\par
悠扬要在神仙境。\par

\section{渔夫舞}
\textbf{史浩}\par\vspace{0.5em}
一引修鳞吾事了。\par
棹船归去歌声杳。\par
门俯清湾山更好。\par
眠到晓。\par
鸣榔艇子方云扰。\par

\section{渔夫舞}
\textbf{史浩}\par\vspace{0.5em}
试倩霜刀登玉缕。\par
银鳞不忍供盘俎。\par
掷向清波方圉圉。\par
休更取。\par
小槽且听真珠雨。\par

\section{渔夫舞}
\textbf{史浩}\par\vspace{0.5em}
莫惜清尊长在手。\par
圣朝化洽民康阜。\par
说与渔家知得否。\par
齐稽首。\par
太平天子无疆寿。\par

\section{望海潮}
\textbf{史浩}\par\vspace{0.5em}
烟笼香径，霞舒花砌，东君绣出芳辰。\par
蝶羽弄轻，鸾声啭巧，嬉嬉舞态歌唇。\par
纶制出严宸。\par
曳耀春品服，荣锡绯银。\par
向此华涂要路，颜色倍精神。\par
珠帘碧甃方新。\par
有兰堂快目，水榭通津。\par
玉斝蘸清，金虬蔼翠，轮蹄尽集簪绅。\par
偕老指双椿。\par
望武林咫尺，同上青云。\par
异日重为此会，应羡凤池人。\par

\section{望海潮}
\textbf{史浩}\par\vspace{0.5em}
烟浓柳径，霞蒸花砌，春深特地芳辰。\par
蝶侣斗狂，莺雏弄巧，嬉嬉舞态歌唇。\par
西铺集簪绅。\par
正桂薰兰玉，天寿松椿。\par
竞捧瑶觥潋滟，来祝纵怀人。\par
当年辍侍严宸。\par
有星轺问俗，熊轼临民。\par
康阜政成，蕃宣治美，归休燕处申申。\par
行庆紫泥新。\par
起钓璜国老，东海之滨。\par
屈指重开此宴，应已拜平津。\par

\section{望海潮}
\textbf{史浩}\par\vspace{0.5em}
熊罴嘉梦，风云享会，磻溪应卜之年。\par
黄菊萃英，红萸酿馥，安排预赏芳筵。\par
环佩拥神仙。\par
向粉额两字，金缕红鲜。\par
最好花裀展处，双凤舞翩翩。\par
人人竞擘香笺。\par
璨珠玑溢目，祝颂无边。\par
彭祖一分，庄椿十倍，千秋未足多言。\par
日驭且停鞭。\par
把燕闲欢乐，分付壶天。\par
笑享亲朋岁岁，春酒庆团圆。\par

\section{感皇恩}
\textbf{史浩}\par\vspace{0.5em}
健卒走红尘，芝封飞到。\par
金缕斜斜印三道。\par
舞鸾翔凤，犹带御炉烟袅。\par
茜衣新象笏，银章好。\par
对此况当，莺花缭绕。\par
画栋翬翬映蓬岛。\par
绣帘初卷，共指松椿偕老。\par
浩歌拚烂醉，金尊倒。\par

\section{感皇恩}
\textbf{史浩}\par\vspace{0.5em}
风雨揽元宵，收灯方了。\par
深院红莲尚围绕。\par
德星同聚，更有祥光临照。\par
始知真洞府，春长好。\par
应是化工，偏怜衰老。\par
剩把青藜作荣耀。\par
正须沈醉，拚却玉山频倒。\par
寄声更漏子，休催晓。\par

\section{满庭芳}
\textbf{史浩}\par\vspace{0.5em}
烘锦花堤，铺绵柳巷，晓来膏雨初晴。\par
画堂初建，碧沼映朱楹。\par
最好芙蓉绣褥，交辉敞、孔雀金屏。\par
那堪更，华裾满坐，和气动欢声。\par
冰清。\par
真美行，棠阴善政，槐市高名。\par
今朝消受得，茜服光荣。\par
况是齐眉并寿，谁云道、乐事难并。\par
相将见，飞凫过阙，除目下彤庭。\par

\section{满庭芳}
\textbf{史浩}\par\vspace{0.5em}
爱日轻融，阴云初敛，一番雪意阑珊。\par
柳摇金缕，梅绽五腮寒。\par
知是东皇翠葆，飞星汉、来止人间。\par
开新宴，笙歌逗晓，和气满城寰。\par
风光，偏舜水，贤侯政美，棠荫多欢。\par
更圜扉草鞠，木索长闲。\par
休向今朝惜醉，红妆映、群玉颓山。\par
相将见，宜春帖子，清夜写金銮。\par

\section{满庭芳}
\textbf{史浩}\par\vspace{0.5em}
微霰疏飘，骄云轻簇，短檠暗淡笼纱。\par
冷禁兰帐，清晓忽飞花。\par
已是平芜步阔，那堪更、折竹如蓑。\par
凭栏处，关心一叶，归兴渺无涯。\par
为瑞，已多少，适从狼子，来自龙沙。\par
赖吾皇神武，薄海为家。\par
尽扫腥风杀气，依然放、红日光华。\par
回头看，山蹊水坞，缟带不随车。\par

\section{满庭芳}
\textbf{史浩}\par\vspace{0.5em}
鲸海波澄，棠阴日永，正宜坐啸雍容。\par
岁丰民乐，无讼到庭中。\par
试数循良自古，龚黄外、谁可追踪。\par
那堪更，恩均髦寿，良会此宵同。\par
璇穹。\par
占瑞处，荧煌五马，璀璨群公。\par
盛笙歌落绮，共引髯翁。\par
祗恐芝泥趣召，双旌展、猎猎飞红。\par
须知道，君王渴见，名久在屏风。\par

\section{满庭芳}
\textbf{史浩}\par\vspace{0.5em}
十载江湖，一朝簪组，宠荣曷称衰容。\par
圣恩不许，归卧旧庐中。\par
慨念东山伴侣，烟霞外、久阔仙踪。\par
今何幸，相逢故里，谈笑一尊同。\par
吾州，真幸会，湖边贺监，海上黄公。\par
胜渭川遗老，绛县仙翁。\par
纵饮何辞烂醉，脸霞转、一笑生红。\par
从今后，婆娑化国，千岁乐皇风。\par

\section{满庭芳}
\textbf{史浩}\par\vspace{0.5em}
油幕初开，騂旄前导，暂归梓里舂容。\par
致身槐揆，功在鼎彝中。\par
自是襟怀绝俗，今犹记、笔砚陈踪。\par
张高会，君恩厚赐，乐与故人同。\par
把麾，鄞水上，相看青眼，谁复如公。\par
况亲陪尊俎，笑接群翁。\par
坐上笙歌屡合，须拚到、晓日酣红。\par
公今去，恩波四海，桃李尽东风。\par

\section{满庭芳}
\textbf{史浩}\par\vspace{0.5em}
玉阙朝回，沙堤烟晓，碧幢光动军容。\par
虎符熊轼，行指七闽中。\par
假道吾乡我里，挥金事、思蹑前踪。\par
倾怀处，萤窗雪案，犹说昔年同。\par
相看，俱老大，襟期道义，不为王公。\par
念儿时聚戏，今已成翁。\par
敢借玉壶美酒，还为寿、金盏翻红。\par
仍频祝，中书二纪，寰海振淳风。\par

\section{满庭芳}
\textbf{史浩}\par\vspace{0.5em}
复拥旌麾，重歌褥裤，满城长自春容。\par
搢绅耆旧，欢溢笑谈中。\par
尽道邦君恺悌，逍遥遂、湖海遐踪。\par
今朝会，公真乐善，屈意与人同。\par
恩勤，东道主，挥金汉傅，怀绶朱公。\par
引群仙环拱，欲寿吾翁。\par
春瓮初澄盎绿，春衫更、轻染香红。\par
持杯愿，归登绛阙，花萼醉春风。\par

\section{满庭芳}
\textbf{史浩}\par\vspace{0.5em}
鹤冷风亭，鸿迷烟渚，晓来雪意填空。\par
酿成佳瑞，端为兆年丰。\par
况有神娲妙手，调和得、云彩皆同。\par
楼台上，铺琼缀玉，随步广寒宫。\par
天公。\par
开地轴，八紘混一，莫辨堤封。\par
又须教、归禽狡兽沈踪。\par
坐见花敷万木，谁知道、春已输工。\par
三杯酒，西湖父老，相与话时雍。\par

\section{庆清朝}
\textbf{史浩}\par\vspace{0.5em}
翠竹茎疏，碧溪流浅，绮窗为尔时开。\par
依稀远岸，才见一点寒梅。\par
冷定半疑是雪，因风还度暗香来。\par
醉清兴，瘦策过桥，黄帽青鞋。\par
繁枝正微雨后，似怨人知晚，泪浥冰腮。\par
殷勤百绕，留连踏遍莓苔。\par
报道玉人睡觉，菱花初试晓妆台。\par
携归去，粉额殢人，比并轻抬。\par

\section{蓦山溪}
\textbf{史浩}\par\vspace{0.5em}
清淡无限，林下逢人少。\par
骑马踏红尘，恁区区、何时是了。\par
名场利海，毕竟白头翁，山簇翠，水拖蓝，只个生涯好。\par
君侯洒落，卜筑开冰沼。\par
三径直危楼，遍岩隈、幽花香草。\par
风勾月引，馀事作诗人，词歌雪，气凌云，寒瘦伦郊岛。\par

\section{青玉案}
\textbf{史浩}\par\vspace{0.5em}
玉姬曾向瑶池舞。\par
轻掷霓裳忤王母。\par
从此烟消飞鹤驭。\par
一来人世，有缘相遇，得得为鸳侣。\par
年年此际霞觞举。\par
彩笔香笺染新句。\par
休饵灵砂奔月去。\par
齐眉不老，直须携手，同上青冥路。\par

\section{青玉案}
\textbf{史浩}\par\vspace{0.5em}
涌金斜转青云路。\par
溯衮衮、红尘去。\par
春色勾牵知几度。\par
月帘风幌，有人应在，唾线馀香处。\par
年来不梦巫山暮。\par
但苦忆、江南断肠句。\par
一笑匆匆何尔许。\par
客情无奈，夜阑归去，簌簌红空雨。\par

\section{青玉案}
\textbf{史浩}\par\vspace{0.5em}
年来减却风情大。\par
百样收心待不作。\par
恰恨仙翁停画舸。\par
雪中把酒，美人频为，浅破樱桃颗。\par
清歌谁许阳春和。\par
悄不放、遥空片云过。\par
惊落梁尘浑可可。\par
一声啭处，故园春近，桃李还知麽。\par

\section{西江月}
\textbf{史浩}\par\vspace{0.5em}
红蓼千堤挺蕊，苍梧一叶辞柯。\par
夜阑清露泻银河。\par
洗出芙蓉半朵。\par
解带初开粉面，绕梁还听珠歌。\par
心期端的在秋波。\par
想得今宵只我。\par

\section{喜迁莺}
\textbf{史浩}\par\vspace{0.5em}
凤阙朱旗展，弄罢五弦，南薰敲竹。\par
雨糁桃蹊，钱浮荷沼，一瞬染成新绿。\par
玉皇香案吏，曾是时、鹤飞江国。\par
对此际，每丹霄效瑞，非烟郁郁。\par
卜筑。\par
陶山曲，风榭月台，图画应难足。\par
绿绮春浓，青蛇星烂，肯便稳栖烟麓。\par
玳筵称寿，清皓齿、霏霏珠玉。\par
竞屈指，看芝封紫检，鸣驺入谷。\par

\section{喜迁莺}
\textbf{史浩}\par\vspace{0.5em}
征鸿回北。\par
正雪洗烧痕，千岩匀绿。\par
鱼纵新漪，梅繁断岸，春到鉴湖一曲。\par
满城绣帘珠幌，暖响聒天丝竹。\par
渐向晚，放芙蕖千顷，交辉华烛。\par
贤牧。\par
棠阴静，康阜政成，褒诏来黄屋。\par
玉笋光寒，紫荷香润，人道此装须趣。\par
且出催花银漏，恣饮宝觥醽醁。\par
向明岁，看传柑归去，腰横金粟。\par

\section{喜迁莺}
\textbf{史浩}\par\vspace{0.5em}
才过元夕。\par
送宴赏未阑，欢娱无极。\par
且莫收灯，仍休止酒，留取凤笙龙笛。\par
金马玉堂学士，当此同开华席。\par
最堪爱，是兰膏光在，金釭连壁。\par
难觅。\par
交欢处，杯吸百川，雅量皆勍敌。\par
老子衰迟，居然怀感，厚意怎生酬得。\par
况已倦游客路，一志归安泉石。\par
但屈指，愿诸贤衮绣，联飞鹏翼。\par

\section{喜迁莺}
\textbf{史浩}\par\vspace{0.5em}
谯门残月。\par
正画角晓寒，梅花吹彻。\par
瑞日烘云，和风解冻，青帝乍临东阙。\par
暖响土牛箫鼓，夹路珠帘高揭。\par
最好是，看彩幡金胜，钗头双结。\par
奇绝。\par
开宴处，珠履玳簪，俎豆争罗列。\par
舞袖翩翻，歌声缥缈，压倒柳腰莺舌。\par
劝我应时纳祜，还把金炉香爇。\par
愿岁岁，这一卮春酒，长陪佳节。\par

\section{喜迁莺}
\textbf{史浩}\par\vspace{0.5em}
雪消春浅。\par
听爆竹送穷，椒花待旦，系马合簪，鸣鸦列炬，几处玳筵开宴。\par
介我百千眉寿，齐捧玉壶金盏。\par
最奇绝，是小桃新坼，争妍粉面。\par
女伴。\par
频告语，守岁通宵，莫放笙歌散。\par
酒晕朝霞，寒欺重翠，却忆凤屏香暖。\par
笑拂满身花影，遥指珠帘深院。\par
待到了，道一声稳睡，明年相见。\par

\section{喜迁莺}
\textbf{史浩}\par\vspace{0.5em}
凭高寓目。\par
爱屹起四窗，云南云北。\par
缥缈烟霞，萧森松竹，多少洞天岩谷。\par
著向十洲三岛，入海何妨登陆。\par
要知处，在皇家新赐，西湖一曲。\par
林麓。\par
真胜概，樊榭鹿亭，百卉生幽馥。\par
绿绮春浓，青蛇星烂，隔断世间尘俗。\par
笑呼羡门俦侣，时引宝觞醽醁。\par
醉和醒，但南山之寿，难忘勤祝。\par

\section{点绛唇}
\textbf{史浩}\par\vspace{0.5em}
我为劳生，自怜浪迹天涯遍。\par
如今春换。\par
又是孤萍断。\par
谁信年时，老子情非浅。\par
思量见。\par
画楼天远。\par
花倚夕阳院。\par

\section{点绛唇}
\textbf{史浩}\par\vspace{0.5em}
千里欢谣，使君美政高三辅。\par
沸天箫鼓。\par
笑拥锋车去。\par
卧辙攀辕，漫拟双旌住。\par
还知否。\par
禁林深处。\par
已辟金闺路。\par

\section{点绛唇}
\textbf{史浩}\par\vspace{0.5em}
翠幄园林，火云方绽南薰起。\par
玉轮天外。\par
夜色凉如水。\par
况有清歌，劝我尊浮蚁。\par
拚沈醉。\par
万花丛里。\par
一枕朦胧睡。\par

\section{点绛唇}
\textbf{史浩}\par\vspace{0.5em}
曾到蟾宫，玉轮乞得长随手。\par
数声轻叩。\par
已有锵琼玖。\par
最好长清，浑不惊秋候。\par
歌阑后。\par
那回辞酒。\par
笑把遮檀口。\par

\section{木兰花慢}
\textbf{史浩}\par\vspace{0.5em}
喜阳和应律，启佳气、满寰瀛。\par
正雪洗疏梅，云浮淡月，昨夜生明。\par
熊罴信占梦好，当年相阀再蟠英。\par
收拾仙风道韵，萃兹一点台星。\par
功名。\par
壮岁逢真。\par
主紫橐、耀西清。\par
向玉笋光中，瑶林宴里，来拥双旌。\par
青毡家旧物，看长参、鼎鼎乐升平。\par
春醑休辞介寿，鹤书已播彤廷。\par

\section{临江仙}
\textbf{史浩}\par\vspace{0.5em}
衮绣蝉联三重客，朝回晓日曈昽。\par
绿杨门巷拥花骢。\par
喜承天上语，来作主人公。\par
况值瑶林风露爽，冰轮碾上晴空。\par
桂香和影堕金锺。\par
莫辞通夕醉，明日是秋中。\par

\section{临江仙}
\textbf{史浩}\par\vspace{0.5em}
槛竹敲风初破睡，楚台梦雨精神。\par
背屏斜映小腰身。\par
山明双翦水，香满一钗云。\par
炉袅金丝帘窣地，绮窗秋静无尘。\par
半钩春笋带湘筠。\par
兰亭初写就，愁杀卫夫人。\par

\section{临江仙}
\textbf{史浩}\par\vspace{0.5em}
绣幕罗裙风冉冉，象床毡幄低垂。\par
兽炉香袅锦屏围。\par
不贪敧钿枕，偏爱倚花枝。\par
软玉红绡轻暖透，温温翠袖扶持。\par
正忺安稳坐移时。\par
雨云忘峡梦，身境是瑶池。\par

\end{document}